% diagrams/firstwords/tick.tex -- el reloj de pie del cuento de Grandma,
% con el ratoncito que vive dentro, asomado por su agujero. Tick, tock!
\begin{tikzpicture}[dibujofw]
  % el reloj: la caja, la cabeza con su arco, la esfera y el péndulo
  \filldraw[fill=white] (-1.35,0) rectangle (1.35,0.5);
  \filldraw[fill=white] (-1.1,0.5) rectangle (1.1,4.3);
  \filldraw[fill=white] (-1.3,4.3) -- (-1.3,6.6) arc[start angle=180, end angle=0, radius=1.3]
    -- (1.3,4.3) -- cycle;
  \filldraw[fill=white] (0,5.9) circle (0.95);
  \foreach \a in {0,30,...,330} {\fill ({0.75*cos(\a)},{5.9+0.75*sin(\a)}) circle (0.06);}
  \draw (0,5.9) -- (0,6.5); \draw (0,5.9) -- (0.42,5.72);
  \filldraw[fill=white] (-0.6,1.6) rectangle (0.6,3.9);
  \draw (0,3.9) -- (0,2.45);
  \filldraw[fill=white] (0,2.2) circle (0.3);
  % el ratón, al lado del reloj
  \begin{scope}[shift={(3.2,0)}, scale=1.5]
    \draw (-0.3,0.3) .. controls (-0.9,0.1) and (-1.1,0.6) .. (-0.8,0.9);
    \filldraw[fill=white] (-0.35,0) .. controls (-0.5,0.6) and (0.1,0.95) .. (0.55,0.55)
      .. controls (0.8,0.4) and (0.95,0.25) .. (1.15,0.2) .. controls (0.9,0.05) .. (0.6,0)
      -- cycle;
    \filldraw[fill=white] (0.3,0.75) circle (0.28);
    \filldraw[fill=white] (0.3,0.75) circle (0.13);
    \fill (0.78,0.4) circle (0.05);
    \fill (1.15,0.2) circle (0.07);
  \end{scope}
  \draw (-2,0) -- (5.4,0);
\end{tikzpicture}
